\documentclass[11pt,reqno]{amsart}
\usepackage[T1]{fontenc}
\usepackage{amsmath,amssymb,amsthm}
\usepackage[margin=1in]{geometry}
\usepackage[colorlinks=true,linkcolor=blue]{hyperref}
\newtheorem{theorem}{Theorem}[section]
\newtheorem{lemma}[theorem]{Lemma}
\newtheorem{proposition}[theorem]{Proposition}
\newtheorem{corollary}[theorem]{Corollary}
\theoremstyle{definition}
\newtheorem{definition}[theorem]{Definition}
\newtheorem{hypothesis}[theorem]{Hypothesis}
\theoremstyle{remark}
\newtheorem{remark}[theorem]{Remark}
\emergencystretch=3em
\title{REV-PLAN: exact shared interfaces for the AKLT gap revision}
\date{}
\begin{document}
\maketitle
\noindent These are mathematical statements for parallel proof writing.
No proof or certificate verification is asserted by this interface file.
All expectations and functionals below are finite sums.
The statements with hypotheses are internal implications; the target
main theorem is unconditional after the computational input theorems
are proved. Labels are shared with the complete inventory in PLAN.md.




\section{Result and conventions}

% Owner: A
\begin{definition}\label{d:model}
For integers \(L,N\ge3\), let \(\Lambda=\Lambda_{L,N}=\mathbb Z_L\times\mathbb Z_N\)
with its simple, unoriented nearest-neighbour edge set \(E_{L,N}\).
The Hilbert space is \(\mathcal H_\Lambda=\bigotimes_{x\in\Lambda}V_4\).
Each site carries the spin-\(2\) space \(V_4\), and \(P_{\{x,y\}}\)
is the orthogonal projection onto total spin \(4\) in \(V_4\otimes V_4\).
Fix one orientation of each edge and set
\[
 q_{xy}=z_{x,1}z_{y,2}-z_{x,2}z_{y,1},\qquad
 B_{L,N}=\prod_{e\in E_{L,N}}q_e,\qquad
 H_{L,N}=\sum_{e\in E_{L,N}}P_e.
\]
Every edge is counted once. The spectral gap of \(H_{L,N}\) means its
smallest positive eigenvalue.
\end{definition}


\section{Finite moment algebra and the response input}

% Owner: C
\begin{definition}\label{d:moment}
Write \(\mathbb S^2=\{a\in\mathbb R^3:|a|=1\}\), put
\(\phi=(1+\sqrt5)/2\), \(r=\sqrt{\phi+2}\), and set
\[
 D=\frac1r\{(0,\pm1,\pm\phi),(\pm1,\pm\phi,0),(\pm\phi,0,\pm1)\},
 \qquad
 \mathcal M(f)=\frac1{12}\sum_{a\in D}f(a).
\]
All indicated signs are independent. Define finite probability laws
\(\lambda_{\mathrm U}=\frac1{12}\sum_{a\in D}\delta_a\) and
\(\lambda_{\mathrm D}=\frac35\lambda_{\mathrm U}+\frac25\delta_0\).
Products of these functionals always mean finite sums.
For plane sites, \(d\) denotes ambient \(\ell^1\) distance and
\(\operatorname{dist}_\infty\) denotes ambient \(\ell^\infty\) distance.
Distances to an empty set are \(+\infty\); a product over an empty site
set has its one empty configuration. For finite probabilities, use
\[
 d_{\rm TV}(p,p')=\tfrac12\sum_z|p(z)-p'(z)|,\qquad
 \operatorname{Var}_\mu f=\mu(|f-\mu f|^2),\qquad
 d_H(n,n')=\sum_x\mathbf1_{\{n_x\ne n'_x\}} .
\]
\end{definition}

% Owner: C
\begin{lemma}\label{l:moments}
The set \(D\) has twelve distinct unit vectors and is antipodal.
For \(i,j,k,l\in\{1,2,3\}\),
\[
 \begin{gathered}
 \mathcal M(1)=1,\qquad \mathcal M(n_i)=\mathcal M(n_in_jn_k)=0,\qquad
 \mathcal M(n_in_j)=\delta_{ij}/3,\\
 \mathcal M(n_in_jn_kn_l)=
 (\delta_{ij}\delta_{kl}+\delta_{ik}\delta_{jl}+\delta_{il}\delta_{jk})/15.
 \end{gathered}
\]
For every polynomial \(p\) of total degree at most four and every
\(R\in\mathrm O(3)\), \(\mathcal M(p\circ R)=\mathcal M(p)\).
The same covariance holds for \(\lambda_{\mathrm D}\); its constant
moment is one and its moments of positive degree at most four are
\(3/5\) times those displayed above.
These assertions hold site by site for polynomials of separate degree
at most four in a finite product of spins.
\end{lemma}

% Owner: C
\begin{definition}\label{d:pinmodel}
For finite \(G\subset\mathbb Z^2\), let
\[
 \mathcal L(G)=\{(x,z):x\in G,\ z\notin G,\ d(x,z)=1\},\qquad
 S(G)=\{x\in G:d(x,\mathbb Z^2\setminus G)\le2\}.
\]
A pin assignment is \(\tau\in(\mathbb S^2)^{\mathcal L(G)}\).
For \(Y\subset G\), write \(D_Y=G\setminus Y\) and define
\[
 Z^\tau_{D_Y}[f]=
 \sum_{n\in D^{D_Y}}12^{-|D_Y|}f(n)
 \prod_{\substack{\{x,y\}\subset D_Y\\d(x,y)=1}}(1+n_x\cdot n_y)
 \prod_{\substack{(x,z)\in\mathcal L(G)\\x\in D_Y}}
       (1+\tau_{xz}\cdot n_x).
\]
The free sum \(Z^0_{D_Y}[f]\) omits all leg factors.
Write \(Z^\bullet_{D_Y}=Z^\bullet_{D_Y}[1]\),
\(\mathbb E^\bullet_{D_Y}f=Z^\bullet_{D_Y}[f]/Z^\bullet_{D_Y}\),
and \(\Psi_{D_Y}(\tau)=Z^\tau_{D_Y}/Z^0_{D_Y}\).
Here and below \(\bullet\) is either \(\tau\) or \(0\).
The empty partition sum is one. There are no new factors on legs into \(Y\).
Different original legs may carry unrelated unit pins, even if their
outside endpoints coincide.
\end{definition}

% Owner: C
\begin{lemma}\label{l:finitefacts}
Consider a finite network with site laws \(\lambda_{\mathrm U}\) or
\(\lambda_{\mathrm D}\), retained distinct-endpoint bonds
\(1+n_x\cdot n_y\), and retained original unit-pin factors
\(1+\tau_\ell\cdot n_x\), with at most four factors incident to each site.
Its pinned and free partition sums are strictly positive, including
the empty network. Non-neighbours are conditionally independent given
the other spins on every conditioning event of positive probability.
If \(f\) is a product of scalar products of spins and the product of
\(f\) with the bond factors has separate degree at most four at every
site, then \(Z^0[f]\ge0\), hence \(\mathbb E^0 f\ge0\).
All normalized spins have norm at most one.
\end{lemma}

% Owner: C
\begin{lemma}\label{l:35}
For \(0\le k\le3\) and \(|a_j|\le1\), the probability law proportional to
\(\prod_{j=1}^k(1+n\cdot a_j)\lambda_{\mathrm U}(\mathrm dn)\)
has mean of norm at most \(3/5\). The same bound holds with
\(\lambda_{\mathrm D}\) in place of \(\lambda_{\mathrm U}\).
Moreover,
\[
 \left\{
 \frac{\mathcal M\!\left(n\prod_{j=1}^3(1+n\cdot a_j)\right)}
      {\mathcal M\!\left(\prod_{j=1}^3(1+n\cdot a_j)\right)}
       :a_1,a_2,a_3\in\mathbb S^2
 \right\}=\{b\in\mathbb R^3:|b|\le3/5\}.
\]
Consequently an unpinned site having at most three present factors
has conditional mean of norm at most \(3/5\), on every positive
conditioning fibre.
\end{lemma}

% Owner: C
\begin{lemma}\label{l:exact}
At each site of a finite set choose any positive weighted unit-vector
law whose moments through degree four equal those of \(\mathcal M\).
A polynomial of separate degree at most four has the same product
functional under these laws as under \(\lambda_{\mathrm U}\) at all
sites. The analogous assertion holds after replacing any chosen site
laws by \(3/5\) times that law plus \(2/5\,\delta_0\).
In particular this applies to partition products of factors
\(1\pm n_x\cdot n_y\) and \(1\pm n_x\cdot a\), \(|a|=1\), whenever
each variable occurs at most four times, and to inserted observables
only when their degrees are included in this same bound.
\end{lemma}

% Owner: C
\begin{definition}[(P)]\label{h:P}
For \(A\ge2\) and \(0<\theta<1\), property \(\mathsf P(A,\theta)\)
means that, for every finite \(G\subset\mathbb Z^2\), every unit-pin
assignment on \(\mathcal L(G)\), every
\(Y\subset G\setminus S(G)\), and all \(u,v\in G\setminus Y\) with
\(d(u,v)=2\) and a common neighbour in \(Y\),
\[
 \left|\mathbb E^\tau_{G\setminus Y}(n_u\cdot n_v)
             -\mathbb E^0_{G\setminus Y}(n_u\cdot n_v)\right|
 \le A\theta^{\min\{d(u,S(G)),d(v,S(G))\}} .
\]
The distances are ambient plane distances to the original \(S(G)\);
removed sites are not added to the boundary.
\end{definition}


\section{The quantum model, uniqueness, and positive frames}

% Owner: A
\begin{definition}\label{d:poly}
Let \(V_d\) be the homogeneous binary polynomials of degree \(d\).
Its monomials \(z_1^{d-j}z_2^j\), \(0\le j\le d\), are orthogonal with
squared norms \(\binom dj^{-1}\). The action is
\((U_gp)(z)=p(g^{-1}z)\). Use \(V_4\) at each physical site.
For each of the finitely many frame letters \(a\) choose one unit
spinor \(\zeta(a)\in\mathbb C^2\) with
\(\zeta(a)\zeta(a)^*=(I+a\cdot\sigma)/2\).
No spinor choice on the whole sphere is part of the definitions.
\end{definition}

% Owner: A
\begin{definition}\label{d:frame}
A frame is a finite set \(A\subset\mathbb S^2\) with positive weights
\(w(a)\) summing to one such that
\(\sum_{a\in A}w(a)p(a)=\mathcal M(p)\)
for every real polynomial \(p\) of degree at most four.
\end{definition}

% Owner: A
\begin{lemma}\label{l:tight}
For every frame \((A,w)\) and every \(p\in V_4\),
\[
 \sum_{a\in A}5w(a)|p(\zeta(a))|^2=\|p\|_{V_4}^2.
\]
For site frames \((D_x,w_x)\), set \(\Omega=\prod_xD_x\),
\(\pi(s)=\prod_x5w_x(s_x)\), and
\(\iota\varphi(s)=\varphi((\zeta(s_x))_x)\).
Then \(\iota:\mathcal H_\Lambda\to\ell^2(\Omega,\pi)\) is an isometry.
\end{lemma}

% Owner: A
\begin{lemma}\label{l:threeframes}
Let
\[
 T=\begin{pmatrix}99/101&-20/101&0\\20/101&99/101&0\\0&0&1\end{pmatrix},
 \quad D^{(j)}=T^jD\ (j=0,1,2),\quad
 b=\frac{5-\sqrt5}{505},\quad r_b=\frac b{2-b}.
\]
The matrix \(T\) is in \(\mathrm{SO}(3)\).
Each \(D^{(j)}\), with uniform weights \(1/12\), is an antipodal frame.
For \(i\ne j\), \(a\in D^{(i)}\), \(a'\in D^{(j)}\),
\[
 0<b<1,\qquad b\le1\pm a\cdot a'\le2-b,\qquad0<r_b<1.
\]
The common lower bound \(b\) is attained, and distinct frames are
disjoint even up to antipodal signs.
\end{lemma}

% Owner: A
\begin{definition}\label{d:framemodel}
Given any proper three-colouring \(\kappa\) of a plane or torus,
put \(D_x=D^{(\kappa(x))}\), with weights \(1/12\).
On the torus define
\[
 \mu^{\mathrm{af}}_{\kappa,L,N}(s)=
 \frac1{Z^{\mathrm{af}}}\prod_x\frac1{12}
       \prod_{\{x,y\}\in E_{L,N}}(1-s_x\cdot s_y).
\]
On the plane, for finite \(W\subset\mathbb Z^2\) and
\(\eta_z\in D_z\) at exterior neighbour sites, define
\[
 \mu^{+,\eta}_{\kappa,W}(n)=\frac1{Z^{+,\eta}_{\kappa,W}}
 \prod_{x\in W}\frac1{12}
 \prod_{\substack{\{x,y\}\subset W\\d(x,y)=1}}(1+n_x\cdot n_y)
 \prod_{\substack{x\in W,z\notin W\\d(x,z)=1}}(1+n_x\cdot\eta_z).
\]
Only the finitely many needed exterior letters are specified.
These laws have full support. Write \(E_U\) for conditioning off \(U\),
\(\operatorname{Var}_U f=E_U|f-E_Uf|^2\), and
\[
 \mathcal D_\mu(f)=\sum_x\mu(\operatorname{Var}_x f),\qquad
 \lambda(\mu)=\inf_{\operatorname{Var}_\mu f>0}
             \frac{\mathcal D_\mu(f)}{\operatorname{Var}_\mu f}.
\]
All generators have one update per physical site. A gap lower bound is
always equivalently a Poincare inequality, including one-point spaces.
\end{definition}

% Owner: A
\begin{lemma}\label{l:planeenergy}
All the frame specifications above are consistent under conditioning
and have nearest-neighbour Markov dependence. Uniformly over proper
colourings, finite volumes, admissible frame-valued boundaries, and
fully specified one-site fibres, every letter has conditional
probability at least
\[
 \alpha=\frac1{12}r_b^4\in(0,1].
\]
The same lower bound holds after conditioning on fewer letters.
Any prescribed assignment on \(k\) sites, under arbitrary consistent
partial conditioning off those sites, has conditional probability
at least \(\alpha^k\).
For a finite plane block \(Q\), the single-site heat bath of every conditional block law in \(Q\) has Poincare gap at least
\(r_b^{|E(Q)|}\), where \(E(Q)\) is its internal edge set.
\end{lemma}


\section{Windows, stopped coupling, and the periodic block gap}

% Owner: A (frozen shared definition)
\begin{definition}\label{d:shells}
For \(a\in\mathbb Z^2\), integer \(u\ge1\), put
\(C_u(a)=a+\{-u,\ldots,u\}^2\).
For \(\operatorname{dist}_\infty(y,C_u(a))=1\), let
\(s=u^{1/4}\), \(k_i=\lfloor is\rfloor\), \(i_1=\lfloor s\rfloor+1\).
For an integer \(l_0\ge1\) define
\[
 \begin{split}
 F_i&=\{x\in C_u(a):|x-y|_\infty>(i-1)s\},\\
 D_i&=\{x\in C_u(a):|x-y|_\infty=k_i\},\\
 A_i&=\{x\in D_i:\operatorname{dist}_\infty
                    (x,\mathbb Z^2\setminus C_u(a))\le l_0\},\qquad
 B_i=D_i\setminus A_i .
 \end{split}
\]
Let \(I_0\) consist of the integers \(2\le i\le i_1\) for which
some point of \(D_i\) is at distance strictly less than \(s/3\)
from a corner site of \(C_u(a)\).
All these are sets of physical sites.
\end{definition}

% Owner: A (contract); B (establishes it)
\begin{definition}[(W)]\label{d:W}
The family of plane frame specifications has \(\mathsf W(l_0,p_0,M_\star)\)
if \(l_0,M_\star\) are positive integers, \(0<p_0\le1\), and,
for every proper plane three-colouring, every translate \(C_u(a)\)
with integer \(u\ge M_\star\), every \(y\) as in Definition~\ref{d:shells},
every integer \(2\le i\le u^{1/4}\) outside \(I_0\), and every two
admissible frame-valued boundary conditions \(\eta,\xi\) for \(F_i\),
\[
 d_{\mathrm{TV}}\!\left(
       \mu^{+,\eta}_{\kappa,F_i,B_i},
       \mu^{+,\xi}_{\kappa,F_i,B_i}\right)\le1-p_0 .
\]
The subscript \(B_i\) denotes a marginal. Boundaries are arbitrary
within their site alphabets; they need not agree anywhere.
\end{definition}


\section{The cut and the site-scale window theorem}

% Owner: B
\begin{lemma}\label{l:shellgeom}
For \(u\ge M_{\mathrm g}(l_0)=\max\{16,6(l_0+1)\}^4\),
\(|I_0|\le1\), and \(D_i\subset F_i\) for \(2\le i\le i_1\).
For \(2\le i\le i_1\) outside \(I_0\), \(|A_i|\le8l_0\).
Every side of \(y+[-k_i,k_i]^2\) parallel to an edge of \(C_u(a)\)
and meeting \(C_u(a)\) is at distance at least \(s/3\) from the
exterior half-plane beyond that edge.
For \(i\ge2\), every point of \(C_u(a)\setminus F_i\)
within \(\ell^\infty\) distance one of \(F_i\) lies in \(D_{i-1}\);
\(y\) is not within that distance of \(F_i\).
\end{lemma}

% Owner: B
\begin{theorem}\label{t:cutMOS}
If \(\mathsf P(A,\theta)\) holds, then the family of plane frame
specifications has \(\mathsf W(64,p_0,M_\star)\), with common
\(p_0\in(0,1]\) and integer \(M_\star\ge1\) depending only on
\(A,\theta\). The constants work for every proper plane
three-colouring and every translate.
Specifically, let
\[
 E=1+4A+10A\theta^{-4}
       \left[48\theta^{14}\left(17+\frac1{1-\theta}\right)+1\right],
 \qquad p_0=e^{-2E}.
\]
There exists an integer \(M_\star\ge1\) for which (W) holds
with this displayed choice of \(p_0\).
\end{theorem}


\section{Walk expansion without scalar trace masses}

% Owner: C
\begin{definition}\label{d:domain}
Fix finite \(G\), \(S=S(G)\), and original unit pins \(\tau\).
A domain assigns statuses \(\mathrm U,\mathrm D,\mathrm X\);
\(\mathrm X\) sites are removed. Each nonremoved site has a kept
subset of its incident bonds inside \(G\). A bond is present if both
ends are nonremoved and both keep it. A \(\mathrm D\) site has at
most two present bonds. Protected sites in \(S\) have status
\(\mathrm U\) and keep every incident bond inside \(G\).
Use \(\lambda_{\mathrm U},\lambda_{\mathrm D}\) at the corresponding
sites and define
\[
 Z_D^\tau[f]=\sum_n f(n)
       \prod_{\text{present }\{x,y\}}(1+n_x\cdot n_y)
       \prod_{\substack{(x,z)\in\mathcal L(G)\\x\in S}}
                   (1+\tau_{xz}\cdot n_x)
       \prod_{x:\,\mathrm{st}(x)\ne\mathrm X}\lambda_{\mathrm{st}(x)}(n_x).
\]
The sum is over the product of the supports \(D\) at \(\mathrm U\)
sites and \(D\cup\{0\}\) at \(\mathrm D\) sites.
Define the free sum, normalized expectations and \(\Psi_D\) as before.
The class \(\mathcal D_{\mathrm{cons}}(G)\) consists of domains in
which every \(\mathrm U\) site keeps all its bonds, every \(\mathrm D\)
site keeps exactly two, and every bond between nonremoved sites is
kept by both ends or by neither.
An auxiliary network has the same data but permits protected
\(\mathrm U\) sites to have arbitrary kept subsets. It retains every
original leg factor. Neither domains nor auxiliary networks carry
a scalar trace mass. Only the networks occurring below, with at most
four factors at each site, are used.
\end{definition}

% [AMENDED]
% Owner: C
\begin{lemma}\label{l:response}
At a \(\mathrm U\) departure \(y\notin S\), write
\(r_{ab}=\mathbb E^\tau_{E'}P_{ab}-\mathbb E^0_{E'}P_{ab}\).
With three present neighbours,
\[
 \log\rho_j^0-\log\rho_j^\tau
       =\alpha_jr_{fg}+\beta_j(r_{fo}+r_{go}),\qquad
 0<\alpha_j\le\bar\alpha=\frac{18094}{100000},\quad
 0<\beta_j\le\bar\beta=\frac{36751}{100000}.
\]
With two present neighbours the response is \(\beta_j'r_{fo}\),
where \(0<\beta_j'\le35510/100000\le\bar\beta\);
with one it is zero. Also \(\rho_j^0\le1\) and
\(\rho_j^\tau\le12/5\).
A \(\mathrm D\) departure has zero logarithmic response.
\end{lemma}

% Owner: C
\begin{definition}\label{d:pattern}
Let \(\mathsf T=\{z:d(z,0)\le2\}\), with passage site \(0\),
arrival site \(\mathsf s=(0,-1)\), other neighbours
\(\mathsf e=(1,0),\mathsf N=(0,1),\mathsf W=(-1,0)\),
four corners \((\pm1,\pm1)\), and tips \((\pm2,0),(0,\pm2)\).
The external slots are the two external bonds at each corner and
the three at each tip, twenty in total.
A pattern chooses an out-neighbour \(o\in\{\mathsf e,\mathsf N,\mathsf W\}\)
and statuses/kept sets with the following rules.
The passage site is \(\mathrm U\) and keeps precisely its three
non-arrival bonds. The arrival site is removed or is \(\mathrm D\)
keeping two of its three non-passage bonds.
Each other neighbour is \(\mathrm U,\mathrm X\), or \(\mathrm D\)
keeping its passage bond and one other bond; \(o\) is not removed.
Each corner is \(\mathrm U,\mathrm X\), or \(\mathrm D\) keeping
two of its four bonds. Every \(\mathrm U\) site other than the
passage site keeps all four bonds.
Between nonremoved non-tip sites, a bond is kept by both ends or
by neither. Tips are covering \(\mathrm U\) sites; their bond to
a neighbour \(z\) is present exactly when \(z\) is nonremoved and
keeps that bond. No consistency condition with a global tip kept
set is imposed.
For \(m\in(\mathbb S^2)^{20}\), include every external slot factor
at every nonremoved corner and tip, including both corner slots
when the corresponding actual external bonds are missing.
Let \(Z_{\rm cur}(\pi;m)\) be the finite product functional of
all present template bonds and these slots, using the status laws.
Let \(Z_{\rm next}(\pi;m)\) change the passage law to
\(\lambda_{\mathrm D}\) and drop its out-bond.
The resulting class is \(\mathcal P_{\rm ext}\); it is straight
when \(o=\mathsf N\), and turning otherwise.
\end{definition}

% [AMENDED]
% Owner: C
\begin{definition}[(H1)]\label{h:H1}
For \(c_{\mathrm s},c_{\mathrm t}\in(0,1]\), (H1) means that
for every \(\pi\in\mathcal P_{\rm ext}\) and every
\(m\in(\mathbb S^2)^{20}\),
\[
 Z_{\rm cur}(\pi;m)\ge c_{\rm type(\pi)}Z_{\rm next}(\pi;m).
\]
In the paper the fixed thresholds are
\(c_{\mathrm s}=961/1000\), \(c_{\mathrm t}=187/200\).
No minimum, minimizing pins, or optimal inversion is asserted.
\end{definition}

% Owner: C
\begin{definition}\label{d:weights}
Put \(x_{\mathrm s}=1/c_{\mathrm s}\), \(x_{\mathrm t}=1/c_{\mathrm t}\).
A lattice trail \(w=(v_0,\ldots,v_n)\) has distinct unoriented
nearest-neighbour bonds. Give it a designated unit incoming direction
\(a\), and forbid \(v_1-v_0=-a\) when \(n\ge1\); the virtual
incoming bond is not itself a used bond of the trail.
For departure \(i\), \(0\le i<n\), give weight \(1/5\) if \(v_i\)
has been departed from earlier, and otherwise weight \(x_{\mathrm s}/3\)
for continuing in the incoming direction and \(x_{\mathrm t}/3\)
for a turn. Let \(\mathsf W(w)\) be the product, with
\(\mathsf W(\varnothing)=1\).
\end{definition}

% [AMENDED]
% Owner: C
\begin{definition}[(H2)]\label{h:H2}
For \(C_V\ge1\) and \(0<q<1\), (H2) means that for every
integer \(n\ge0\), every start and every designated incoming direction,
\[
 \sum_{\substack{|w|=n\\\text{initial reversal forbidden}}}
                      \mathsf W(w)\le C_Vq^n .
\]
The sum is over plane lattice trails with those specified initial data.
\end{definition}


\section{Pair closure, a positive-vector witness, and renewal}

% [AMENDED]
% Owner: C
\begin{lemma}\label{l:log-mean-bounds}
Let \(\mathcal L(x,y)=(x-y)/(\log x-\log y)\) for \(x\ne y\)
and \(\mathcal L(x,x)=x\), for positive \(x,y\).
Put \(f(s)=1/\mathcal L(s,1)\), so \(f(1)=1\).
For \(3/5\le c\le1\), the function
\[
 G_c(x,y)=\frac1{3\mathcal L(cx,y)}-\frac1{5\mathcal L(x,y)}
\]
is nonincreasing in each positive variable. In particular,
for \(x\ge a=116/125\), \(y\ge1\),
\[
 G_c(x,y)\le f(ac)/3-f(a)/5.
\]
For rational \(0<s<1\), put \(t=(1-s)/(1+s)\) and
\[
 f_-(s)=\frac2{1-s}\sum_{k=0}^{7}\frac{t^{2k+1}}{2k+1},
 \qquad
 f_+(s)=f_-(s)+\frac{2t^{17}}{17(1-t^2)(1-s)}.
\]
Then \(f_-(s)\le f(s)\le f_+(s)\).
For the two rows of the following table, all numbers are exact rationals:
\[
 \begin{array}{c|c|ccc}
 \tau&c_\tau&a_\tau&b_\tau&b'_\tau\\ \hline
 {\rm s}&961/1000&145217/10^6&352782/10^6&340008/10^6\\
 {\rm t}&187/200&150000/10^6&357549/10^6&344661/10^6
 \end{array}
\]
They satisfy
\[
 \frac{f_+(ac_\tau)}3-\frac{f_-(a)}5\le a_\tau,\qquad
 \frac{f_+(ac_\tau)}3\le b_\tau,\qquad
 \frac{f_+(c_\tau)}3\le b'_\tau,\qquad
 0<b'_\tau\le b_\tau .
\]
\end{lemma}

% [AMENDED]
% Owner: C
\begin{lemma}\label{l:response-typed}
Fix \(c_{\rm s}=961/1000\), \(c_{\rm t}=187/200\) and assume (H1).
At a covered \(\mathrm U\) departure of type \(\tau\in\{{\rm s},{\rm t}\}\),
the response coefficients in Lemma~\ref{l:response} satisfy
\[
 0<\alpha_j\le a_\tau,\qquad 0<\beta_j\le b_\tau
       \quad\hbox{with three present neighbours},\qquad
 0<\beta'_j\le b'_\tau\quad\hbox{with two}.
\]
Here the exact constants are those of Lemma~\ref{l:log-mean-bounds};
in particular \(a_{\rm t}=150000/10^6\).
With \(a=116/125\), their analytic bounds are
\[
 \beta_j\le f(ac_\tau)/3,\qquad
 \beta'_j\le f(c_\tau)/3,\qquad
 \alpha_j\le f(ac_\tau)/3-f(a)/5.
\]
These follow using \(I^\tau\ge c_\tau J^\tau\),
\(I^0\ge J^0\ge1\), and \(J^\tau\ge a\).
Also \(a_\tau\le18094/100000\) and
\(b_\tau,b'_\tau\le36751/100000\).
\end{lemma}

% [AMENDED]
% Owner: C
\begin{definition}\label{d:kernel}
For \(X=B,C\), fix \(c=0\), \(u=(-1,0)\), and \(v=(1,0)\)
for \(B\), \(v=(0,1)\) for \(C\).
Let \(\mathcal T_n^X\) be the length-\(n\) plane trails avoiding \(c\),
from \(u\) to their first hit of \(v\), with incoming direction \(u-c\).
A departure is a lattice revisit when its site was departed from earlier.
At a lattice first visit \(i\), put \(x_i=x_{\tau_i}\) according to
its straight/turn type. Of the three neighbours other than the arrival
site, write \(o\) for the outgoing neighbour and \(f,g\) for the others.
A neighbour is dead if it is \(c\) or was departed from twice before \(i\).
The outgoing neighbour of a true trail is not dead.
If \(f,g\) are both live, assign coefficient \(a_{\tau_i}\) to
\(\{f,g\}\), and \(b_{\tau_i}\) to each of \(\{f,o\},\{g,o\}\).
If exactly one of \(f,g\) is live, assign \(b'_{\tau_i}\) to its
pair with \(o\); if neither is live there is no response pair.
The constants are those of Lemma~\ref{l:log-mean-bounds}.
A live pair \(p\) has type \(Y=B,C\) and shift
\[
 s(p)=\max_{z\in p}\min\{d(z,u),d(z,v)\}.
\]
Let \(\kappa_{Y,s}(i)\) sum its assigned coefficients.
For any fixed, \(\theta\)-independent orderings \(\sigma_w\) of lattice
first visits, define
\[
 \begin{split}
 M_{XY,s}
 &=\sum_{n\ge1}\sum_{w\in\mathcal T_n^X}
     3^{-n}(3/5)^{\operatorname{rev}(w)}
     \sum_k\kappa_{Y,s}(\sigma_w(k))
                   \prod_{l\le k}x_{\sigma_w(l)},\\
 M_\theta[X,Y]&=\sum_{s\ge0}M_{XY,s}\theta^{-s}.
 \end{split}
\]
These are nonnegative series. The identity
\(3^{-n}(3/5)^r=3^{-(n-r)}5^{-r}\) retains the mass-free
trace coefficients. This \(M\) is the modified type-dependent kernel,
not the original uniform-coefficient matrix.
\end{definition}

% [AMENDED]
% Owner: C
\begin{definition}\label{d:kernel-order}
Fix
\[
 h=(1,h_C),\quad h_C=\frac{5404319552844595}{4503599627370496},
 \qquad \rho_0=\frac{993015}{10^6},\qquad L=12 .
\]
For a first visit set \(\kappa_h(i)=\sum_{Y,s}\kappa_{Y,s}(i)h_Y\).
For a closed trail of length \(n\le L\), let \(\sigma_w\) order its
first visits by increasing
\((x_i-1)/(x_i\kappa_h(i))\), with value \(+\infty\) when
\(\kappa_h(i)=0\), breaking ties by increasing path index.
For \(n>L\), list straight first visits in path order and then
turning first visits in path order.
Hereafter \(M_{XY,s}\) and \(M_\theta\) use exactly these orderings.
They do not depend on \(\theta\).
\end{definition}

% [AMENDED]
% Owner: C
\begin{definition}[(H3\(_h\))]\label{h:H3}
For the type-dependent kernel and fixed orderings of
Definitions~\ref{d:kernel}--\ref{d:kernel-order}, (H3\(_h\)) is the
following mathematical assertion, with the specified \(h,\rho_0\):
for \(X=B,C\) and every integer \(R_0\ge0\),
\[
 \sum_{n=1}^{R_0}\sum_{w\in\mathcal T_n^X}
 3^{-n}(3/5)^{\operatorname{rev}(w)}
 \sum_k\kappa_h(\sigma_w(k))\prod_{l\le k}x_{\sigma_w(l)}
       \le\rho_0h_X .
\]
Equivalently, the full nonnegative sums satisfy \(M_1h\le\rho_0h\).
No bound for the former uniform-coefficient matrix is part of this
contract. Its proof in this revision uses the stationary lattice
Bellman certificates, with no periodic spatial walk or remainder.
\end{definition}

% [AMENDED]
% Owner: C
\begin{lemma}\label{l:Mtheta}
Assume (H2) and (H3\(_h\)) for the type-dependent kernel.
For every \(q^2<\theta\le1\), all entries of \(M_\theta\) are
finite. There exist \(q^2<\theta<1\) and \(\rho_0<\rho<1\)
such that the same vector \(h\) satisfies \(M_\theta h\le\rho h\).
Every shift of a length-\(n\) trail is at most \(n/2+1\), and
each matrix entry is bounded by
\[
 \sum_{n\ge1} n(\bar\alpha+2\bar\beta)
                      C_V q^n\theta^{-n/2-1}.
\]
\end{lemma}

% [AMENDED]
% Owner: C
\begin{theorem}\label{t:rec}
Fix \((c_{\rm s},c_{\rm t})=(961/1000,187/200)\).
Assume (H1), (H2) with any \(C_V\ge1\), \(0<q<1\), and
(H3\(_h\)) for the type-dependent kernel and fixed orderings above.
Then there exist \(A\ge2\) and \(0<\theta<1\) such that
\(\mathsf P(A,\theta)\) holds.
One may take the \(\theta,\rho,h\) supplied by the shifted-witness
lemma and
\[
 A=\frac{\max h}{\min h}\,
       \frac{\max\{2,\,2(12/5)^{24}(C_V/(1-q))^2\}}{1-\rho}.
\]
\end{theorem}


\section{Exact statements of the computational inputs}

% [AMENDED]
% Owner: C
\begin{theorem}\label{t:H1}
For every \(\pi\in\mathcal P_{\rm ext}\) and every
\(m\in(\mathbb S^2)^{20}\),
\[
 Z_{\rm cur}(\pi;m)\ge
 \begin{cases}
 (961/1000)Z_{\rm next}(\pi;m),&\pi\text{ straight},\\
 (187/200)Z_{\rm next}(\pi;m),&\pi\text{ turning}.
 \end{cases}
\]
\end{theorem}

% [AMENDED]
% Owner: C (theorem text and proof account); separate certificate owner supplies H2/H3 witnesses
\begin{theorem}\label{t:inputs}
At \(c_{\rm s}=961/1000\), \(c_{\rm t}=187/200\), (H1) holds
by the operator-only Gram/residual certificates of Appendix D.
Property (H2) holds with
\[
 q=\frac{1119948287353}{1120437614225}<\frac{2499}{2500}<1,
 \qquad C_V=10 .
\]
For the type-dependent coefficients of Lemma~\ref{l:log-mean-bounds},
the orderings and vector \(h\) of Definition~\ref{d:kernel-order},
and \(\rho_0=993015/10^6\), property (H3\(_h\)) holds.
The latter assertion is supplied by five stationary lattice Bellman
certificates with memory \(m=10\), exact prefix \(L=12\ge m+2\),
radius \(R=70\), value scale \(2^{16}\), and coefficient scale
\(2^{48}\), using the checked 3741-state automaton quotient.
\end{theorem}


\section{The small H2 test and stationary H3 certificates}

% [AMENDED]
% Owner: C
\begin{definition}\label{d:memory-automaton}
Let \(\mathcal E=\{{\rm E}=(1,0),{\rm N}=(0,1),
{\rm W}=(-1,0),{\rm S}=(0,-1)\}\).
For \(m\ge1\), \(\Omega_m\) consists of sequences
\(\omega=(d_1,\ldots,d_m)\in\mathcal E^m\) whose path from zero
uses distinct unoriented bonds. Write \(z_0=0\),
\(z_j=\sum_{i\le j}d_i\), and \(a(\omega)=d_m\).
A direction \(d\) is admissible if \(d\ne-d_m\) and
\(\{z_m,z_m+d\}\) is not a used bond of this window.
Its successor is \(\omega d=(d_2,\ldots,d_m,d)\).
Its kind \(k(\omega,d)\) is \({\rm r}\) if
\(z_m\in\{z_0,\ldots,z_{m-1}\}\), otherwise \({\rm s}\)
if \(d=d_m\), and \({\rm t}\) otherwise.
Direction labels remain absolute lattice directions.
\end{definition}

% [AMENDED]
% Owner: C
\begin{definition}\label{d:memorygrowth}
On \(\Omega_m\), define the nonnegative one-step matrix
\[
 \mathsf A_m(\omega,\omega')=
 \sum_{\substack{d\ {\rm admissible}\\\omega d=\omega'}}
 \gamma(\omega,d),\qquad
 \gamma(\omega,d)=
 \begin{cases}1/5,&k={\rm r},\\
 x_{\rm s}/3,&k={\rm s},\\x_{\rm t}/3,&k={\rm t}.
 \end{cases}
\]
For a prescribed start and incoming direction with initial reversal
forbidden, let \(A_n(\omega)\) be the true total \(\mathsf W\)-weight
of length-\(n\) trails whose last \(m\) steps form \(\omega\),
for \(n\ge m\), and let \(W_n=\sum_{|w|=n}\mathsf W(w)\).
The adopted finite test is \(m=6\) and extension length \(1\).
\end{definition}

% [AMENDED]
% Owner: C
\begin{lemma}\label{l:alm}
For \(n\ge m\),
\[
 A_{n+1}(\omega')\le\sum_\omega A_n(\omega)\mathsf A_m(\omega,\omega').
\]
If \(v_\omega>0\), \(0<q<1\), and \(\mathsf A_mv\le qv\),
then (H2) holds with
\[
 C_V=\max\left\{
      \max_{0\le r<m}W_rq^{-r},\
      \frac{\max v}{\min v}W_mq^{-m}\right\}.
\]
Any larger finite constant is also valid.
\end{lemma}

% [AMENDED]
% Owner: C
\begin{lemma}\label{l:H2-dihedral}
The dihedral group of the square acts on \(\Omega_m\) by applying
the same lattice symmetry to every direction. This action preserves
admissibility, kind, successor and the weights of \(\mathsf A_m\);
thus \(\mathsf A_m(g\omega,g\omega')=\mathsf A_m(\omega,\omega')\).
For any positive vector constant on these orbits, a row inequality
\((\mathsf A_mv)_\omega\le qv_\omega\) checked at one representative
of every orbit holds on every raw window.
For \(m=6\), \(|\Omega_6|=916\) and there are \(115\) orbits.
\end{lemma}

% [AMENDED]
% Owner: C
\begin{proposition}\label{p:H2-test}
At the two fixed H1 thresholds there exists a positive integer
vector \(v\) on \(\Omega_6\), constant on dihedral orbits, with
\[
 \frac{\max v}{\min v}=\frac{250000000000}{46594269069}<6,\qquad
 \max_\omega\frac{(\mathsf A_6v)_\omega}{v_\omega}
       =q_*:=\frac{1119948287353}{1120437614225}<\frac{2499}{2500}.
\]
The weights have common denominator \(D_2=2695605\), with
numerators \(935000,961000,539121\) for \({\rm s},{\rm t},{\rm r}\).
For \(q_*=q_{\rm num}/q_{\rm den}\), the required row test is
\[
 q_{\rm den}\sum_{d\ {\rm admissible}}n_{k(\omega,d)}v_{\omega d}
          \le q_{\rm num}D_2v_\omega .
\]
It suffices to check the \(115\) orbit rows after proving
Lemma~\ref{l:H2-dihedral}; checking all \(916\) raw rows is also valid.
\end{proposition}

% [AMENDED]
% Owner: C
\begin{lemma}\label{l:H2-CV}
With \(m=6\), \(q=q_*\), and the vector in Proposition~\ref{p:H2-test},
the constant supplied by Lemma~\ref{l:alm} is at most \(10\).
Indeed, uniformly over the prescribed start and incoming direction,
\[
 W_r\le x_{\rm t}^r\quad(0\le r\le6),\qquad
 \max\left\{\max_{r<6}W_rq^{-r},\
                 \frac{\max v}{\min v}W_6q^{-6}\right\}
 \le6\left(\frac{200}{187q}\right)^6<10.
\]
Consequently (H2) holds with the explicit rational constants
\(q=q_*\) and \(C_V=10\).
\end{lemma}

% [AMENDED]
% Owner: C
\begin{lemma}\label{l:automaton-quotient}
Let \(\pi:\Omega_m\to Q\) be a surjection onto a finite set.
Suppose that windows with the same image have, for each absolute
direction, identical admissibility and kind, and their admissible
successors have the same image. Suppose also that two such windows
have the same final direction whenever their class has an admissible
non-revisit transition.
Then admissibility, kind and successor define a well-defined automaton
on \(Q\). Every phase-and-position transition whose coefficients
depend only on the absolute direction, kind, position, and final
direction on non-revisit edges factors through this automaton.
For every function \(G(\varphi,q,x)\), its pullback by \(\pi\)
intertwines the raw and quotient backward operators. This
intertwining also holds for each finite stopped continuation sum.
\end{lemma}

% [AMENDED]
% Owner: C
\begin{proposition}\label{p:automaton10}
There are \(60092\) raw windows in \(\Omega_{10}\).
The supplied map to \(Q=\{0,\ldots,3740\}\), with its supplied
representative per class, is surjective and satisfies every
hypothesis of Lemma~\ref{l:automaton-quotient}.
Its admissible labelled edges consist of \(2516\) straight,
\(5032\) turning and \(1872\) revisit edges.
\end{proposition}

% [AMENDED]
% Owner: C
\begin{definition}\label{d:phaseweights}
The phases are \({\rm A},{\rm B},{\rm C},{\rm D}\). Their
diagonal transition weights, in kind order \({\rm s},{\rm t},{\rm r}\), are
\[
 \begin{array}{c|ccc}
 {\rm A}&x_{\rm s}/3&x_{\rm t}/3&1/5\\
 {\rm B}&x_{\rm s}/3&1/3&1/5\\
 {\rm C}&x_{\rm s}/3&1/3&1/5\\
 {\rm D}&1/3&1/3&1/5
 \end{array}
\]
At a first visit with reward \(\kappa_h\), additionally allow
\({\rm A}\to{\rm C}\) with weight \(x_{\rm t}\kappa_h/3\) on a turn,
or \({\rm B}\to{\rm D}\) with weight \(x_{\rm s}\kappa_h/3\)
on a straight passage. No other phase change is allowed.
At a revisit there is no phase change.
Start with weight one in each of \({\rm A},{\rm B}\);
the terminal reward is \(g=(0,0,1,1)\).
For a true trail, use its actual revisit kinds and its actual
type-dependent live-pair reward from Definition~\ref{d:kernel}.
\end{definition}

% [AMENDED]
% Owner: C
\begin{lemma}\label{l:phase-expansion}
For a true closed trail \(w\) ordered with straight first visits
before turning first visits, its \(h\)-weighted contribution to
\(M_1\) equals the sum over phase histories starting in
\({\rm A}\) or \({\rm B}\) and ending in \({\rm C}\) or \({\rm D}\)
of the products of the weights in Definition~\ref{d:phaseweights}.
Equivalently this sum is
\[
 3^{-|w|}(3/5)^{\operatorname{rev}(w)}
       \sum_k\kappa_h(\sigma_w(k))\prod_{l\le k}x_{\sigma_w(l)}.
\]
Splitting such a phase history at any prefix length gives the
product of its exact prefix weight and its corresponding suffix weight.
\end{lemma}

% [AMENDED]
% Owner: C
\begin{definition}\label{d:stationary-walk}
Fix \(X=B,C\), \(d_0\in\{{\rm N},{\rm W},{\rm S}\}\),
the sites \(u,c,v\) of Definition~\ref{d:kernel}, and \(a=u+d_0\).
The state space is
\(\{{\rm A},{\rm B},{\rm C},{\rm D}\}\times Q\times\mathbb Z^2\).
At positions \(c,v\) there are no outgoing transitions.
At any other position \(x\), an admissible automaton direction \(d\)
goes to position \(x+d\) and the quotient successor window.
At \(x=u\) or \(a\), forbid reuse of \(\{u,a\}\) in either
direction and kill every memory-detected revisit departure;
every other allowed departure there has weight \(1/5\),
in the same phase and with no response.
At other positions use the automaton kind and the phase weights.
For a non-revisit response, declare only \(c\) dead among the three
neighbours other than the arrival site, and use the type-dependent
three-neighbour or two-neighbour reward rule of Definition~\ref{d:kernel}.
A transition into \(c\) has zero contribution.
Let \(T_{X,d_0}\) be this nonnegative backward operator, with zero
terminal rows, and \(g(\varphi,q,x)=\mathbf1_{\{x=v\}}g_\varphi\).
The free operator \(\mathsf W\) omits positions and special-site
rules, uses the same quotient automaton and phase weights, and
treats all three response neighbours as live.
\end{definition}

% [AMENDED]
% Owner: C
\begin{lemma}\label{l:stationary-domination}
Let \(P=(v_0,\ldots,v_L)\) be a true prefix avoiding \(c,v\),
with \(v_0=u\), \(v_1=u+d_0\), and \(L\ge m+2\), \(m=10\).
Every true continuation to its first hit of \(v\), together with
each phase history, is dominated by the corresponding path of
\(T_{X,d_0}\) from its prefix phase, last-window class and endpoint.
The comparison holds for every finite suffix horizon and hence
for their nonnegative suprema.
In this comparison:
forgotten true revisits have weight \(1/5\) no greater than any
relaxed first-visit diagonal weight; forgetting dead non-outgoing
neighbours only increases the reward, using \(b'_\tau\le b_\tau\);
and the known-site killing rules cannot kill a true continuation.
Indeed, the known initial departures at times \(0,1\) are outside
the memory at every suffix time, and a memory-detected additional
departure at \(u\) or \(u+d_0\) would be a forbidden third departure.
\end{lemma}

% [AMENDED]
% Owner: C
\begin{lemma}\label{l:Bellman-soundness}
Let \(T\) be a nonnegative, finitely branching weighted transition
operator on a set of states, with zero rows on a terminal set.
Let \(g\ge0\) be supported on that terminal set.
If a finite-valued function \(F\ge0\) satisfies
\(g+TF\le F\), then for every integer \(n\ge0\),
\[
 \sum_{j=0}^{n}T^jg\le F .
\]
Consequently the nonnegative sum of all finite first-hit path
weights times their terminal rewards is finite and at most \(F\).
It suffices that \(F=g\) on terminal states and \(TF\le F\)
elsewhere.
\end{lemma}

% [AMENDED]
% Owner: C
\begin{definition}\label{d:stationary-certificate}
For a case \((X,d_0)\), fix \(R=70\),
\(B_R=\{-R,\ldots,R\}^2\), \(V=2^{16}\), and \(D_3=2^{48}\).
Supply integers \(Z_{\varphi,q}>0\) and
\(U_{\varphi,q,x}\ge0\), \(x\in B_R\).
Set \(\widetilde U_{\varphi,q,x}=U_{\varphi,q,x}\) inside \(B_R\)
and \(Z_{\varphi,q}\) outside, and \(F=\widetilde U/V\), \(z=Z/V\).
For each individual transition coefficient \(t\) of the exact
stationary or free operator, use the proved rational ceiling
\(\widehat t=p(t)/D_3\), \(p(t)=\lceil D_3t\rceil\);
zero coefficients may be omitted.
The certificate checks all the following:
\[
 \begin{gathered}
 \sum_{\text{free edges }e\text{ from }(\varphi,q)}
       p(t_e)Z_{\operatorname{dest}(e)}
       \le D_3Z_{\varphi,q},\qquad Z_{{\rm C},q},Z_{{\rm D},q}\ge V,\\
 0\le U_{\varphi,q,x}\le Z_{\varphi,q},\qquad
 U_{\varphi,q,c}=0,\qquad U_{\varphi,q,v}=Vg_\varphi,\\
 \sum_{\text{stationary edges }e\text{ from }(\varphi,q,x)}
       p(t_e)\widetilde U_{\operatorname{dest}(e)}
       \le D_3 U_{\varphi,q,x}
       \quad(x\in B_R\setminus\{c,v\}).
 \end{gathered}
\]
All tests are integer inequalities on the regenerated automaton,
directions, phases and lattice sites, not on a trusted supplied matrix.
\end{definition}

% [AMENDED]
% Owner: C
\begin{lemma}\label{l:stationary-local-global}
Every certificate of Definition~\ref{d:stationary-certificate}
defines a finite-valued \(F\) on the entire lattice with
\(F=z\) outside \(B_R\), \(F\le z\), and
\[
 g+T_{X,d_0}F\le F .
\]
For an inside site this follows from its checked integer row and
coefficient ceilings. For an outside site every successor value
is at most its corresponding \(z\), and the exact stationary row
is dominated by the free row; the checked free inequality
implies \(\mathsf Wz\le z\).
Thus \(F\) bounds every stopped continuation as in
Lemma~\ref{l:Bellman-soundness}.
\end{lemma}

% [AMENDED]
% Owner: C
\begin{definition}\label{d:prefix-bound}
For each case \((X,d_0)\), let \(S_{X,d_0}\) be the exact
\(h\)-weighted kernel sum over closed trails of lengths \(1\le n\le L=12\)
with first direction \(d_0\), using their prescribed short-trail orderings.
Let \(\mathcal P_{X,d_0,L}\) be all true length-\(L\) prefixes
avoiding \(c\) and \(v\), with first direction \(d_0\).
For \(P\) in this set, let \(p_\varphi(P)\) be the sum of its
exact phase-history products ending in phase \(\varphi\), starting
with weights one in A and B. Use actual trail revisits, actual dead
neighbours and the type-dependent rewards, with the long-trail
straight-first phase weights.
Write \(x(P)\) for its endpoint and \(q(P)\) for its last-window class.
For the potential of its case define the rational number
\[
 J_{X,d_0}=S_{X,d_0}
       +\sum_{P\in\mathcal P_{X,d_0,12}}\sum_\varphi
                    p_\varphi(P)F(\varphi,q(P),x(P)).
\]
A hit at time \(12\) is in \(S_{X,d_0}\), not in the prefix sum.
\end{definition}

% [AMENDED]
% Owner: C
\begin{lemma}\label{l:prefix-domination}
For every case whose integer certificate passes, the \(h\)-weighted
contribution of all true closed trails of that first direction,
over every length, is at most \(J_{X,d_0}\).
This holds uniformly for every finite length cutoff, with the same
upper bound. The short part uses the exact order through \(12\);
all longer trails split uniquely into a surviving length-\(12\)
prefix and a dominated stopped suffix.
\end{lemma}

% [AMENDED]
% Owner: C
\begin{lemma}\label{l:kernel-reflection}
For row B, reflection \((x_1,x_2)\mapsto(x_1,-x_2)\) fixes
\(c,u,v\), interchanges the first directions N and S, and preserves
every trail weight, type, dead pair, shift, \(h\)-reward and
prescribed ordering. Hence the true B/N and B/S kernel contributions
are equal at every length and every finite cutoff.
No corresponding reduction is asserted for row C.
\end{lemma}

% [AMENDED]
% Owner: C
\begin{proposition}\label{p:stationary-check}
For the five cases B/N, B/W, C/N, C/W, C/S, there exist
integer arrays \(Z,U\) satisfying every check in
Definition~\ref{d:stationary-certificate}, with
\(m=10,L=12,R=70\) and the specified scales.
Their regenerated quotient has \(3741\) states and the spatial
square has \(19881\) sites. Their exact rational prefix evaluations obey
\[
 \begin{array}{c|ccccc}
 (X,d_0)&(B,{\rm N})&(B,{\rm W})&(C,{\rm N})&(C,{\rm W})&(C,{\rm S})\\ \hline
 10^9J_{X,d_0}\ \le&
 351207554&290599163&527872005&380754375&273048088
 \end{array}
\]
Each displayed upper bound denotes an exact integer inequality
after clearing the rational prefix denominator.
\end{proposition}

% [AMENDED]
% Owner: C
\begin{lemma}\label{l:stationary-assembly}
With the five certified potentials and their exact prefix bounds,
\[
 (M_1h)_B\le2J_{B,{\rm N}}+J_{B,{\rm W}}\le\rho_0h_B,\qquad
 (M_1h)_C\le J_{C,{\rm N}}+J_{C,{\rm W}}+J_{C,{\rm S}}
                                      \le\rho_0h_C ,
\]
where \(h_C=5404319552844595/4503599627370496\) and
\(\rho_0=993015/10^6<1\).
The same inequalities hold for every finite true-trail cutoff;
therefore (H3\(_h\)) holds for the modified kernel.
\end{lemma}


\section{Operator-only verification of (H1)}

% [AMENDED]
% Owner: C
\begin{lemma}\label{l:Iw}
Let \(Z_w\) be the finite template functional after deleting
the passage site and its incident factors, retaining all other
present bonds and all covered slot factors. Put \(P_{ab}=n_a\cdot n_b\)
for nonremoved sites and zero if either is removed.
For the out-neighbour \(o\) and the other two neighbours \(f,g\),
\[
 Z_{\rm cur}=Z_w[I],\qquad Z_{\rm next}=Z_w[J]>0,\qquad
 I=1+\tfrac13(P_{fg}+P_{fo}+P_{go}),\quad J=1+\tfrac15P_{fg}.
\]
For \(0<c<1\), \(F_c=Z_{\rm cur}-cZ_{\rm next}=Z_w[I-cJ]\), with
\[
 I-cJ=(1-c)+\tfrac13(P_{fo}+P_{go})+(\tfrac13-\tfrac c5)P_{fg}.
\]
Both partition polynomials are affine in each slot vector and
invariant under simultaneous orthogonal transformations of the slots.
\end{lemma}

% [AMENDED]
% Owner: C
\begin{lemma}\label{l:leaf}
In a covered template, integrating a nonremoved tip or corner
with no present internal bond gives a positive factor depending
only on its own slots.
A tip with its one internal bond to \(x\) gives
\(A(1+(3/5)u\cdot n_x)\), \(A>0,\ |u|\le1\).
A corner with exactly one internal bond to \(x\) gives
\(A(1+\rho u\cdot n_x)\), \(A>0,\ |u|\le1\), where
\(\rho=1/2\) for status \(\mathrm U\) and \(\rho=1/3\)
for status \(\mathrm D\).
\end{lemma}

% [AMENDED]
% Owner: C
\begin{definition}\label{d:reduction}
Delete tips and corners with no present internal bond, dropping
their common positive factors. Replace every bonded tip by a
single vector and radius \(3/5\).
In the reduced representation also replace each one-bond corner
by a single vector of radius \(1/2\) or \(1/3\) according to
Lemma~\ref{l:leaf}; in the unreduced representation retain its spin
and two slots. Other bonded corners retain their two slots.
Let \(u=(u_1,\ldots,u_k)\) list these remaining vectors, and define
\(\widehat Z_{\rm cur},\widehat Z_{\rm next}\) by these replacements,
\(\widehat F_c=\widehat Z_{\rm cur}-c\widehat Z_{\rm next}\).
A group is one single vector or the two slots of a retained corner.
\end{definition}

% [AMENDED]
% Owner: C
\begin{corollary}\label{c:reduce}
In either representation, \(\widehat F_c\) is real, affine in
each vector and invariant under simultaneous \(\mathrm O(3)\).
If \(\widehat F_c\ge0\) on \((\mathbb S^2)^k\), then
\(F_c\ge0\) on \((\mathbb S^2)^{20}\).
For each original slot assignment there are a common \(K>0\)
and vectors \(\widehat u\) in the closed unit ball such that
\(Z_{\rm cur}=K\widehat Z_{\rm cur}(\widehat u)\) and
\(Z_{\rm next}=K\widehat Z_{\rm next}(\widehat u)\).
\end{corollary}

% [AMENDED]
% Owner: C
\begin{definition}\label{d:Pauli}
For a real polynomial affine in each vector, expand uniquely
\[
 F(u)=\sum_{a\in\{0,1,2,3\}^k}f_a\prod_i u_i^{a_i},
 \qquad u_i^0=1,\qquad
 \mathcal O(F)=\sum_a f_a\bigotimes_i\sigma_{a_i},\quad\sigma_0=I.
\]
For \(\mathcal I\subset\{1,\ldots,k\}\), invert those vectors,
writing \(u^{\mathcal I}\) and \(F^{\mathcal I}(u)=F(u^{\mathcal I})\).
For \(|a|\le1\), put \(\rho(a)=(I+a\cdot\sigma)/2\).
Set \(S_z=\frac12\sum_i\sigma_3^{(i)}\) and
\(S_\pm=\sum_i(\sigma_1^{(i)}\pm\mathrm i\sigma_2^{(i)})/2\).
\end{definition}

% [AMENDED]
% Owner: C
\begin{lemma}\label{l:inversion}
For \(|u_i|\le1\), each \(\rho(u_i)\) is positive semidefinite
of trace one, and
\[
 F(u)=\operatorname{Tr}\!\left[
      \mathcal O(F^{\mathcal I})
                 \bigotimes_i\rho((u^{\mathcal I})_i)\right].
\]
Thus a single inversion \(\mathcal I\) with
\(\mathcal O(F^{\mathcal I})\ge0\) implies \(F\ge0\)
on the product of closed unit balls, in particular on unit vectors.
\end{lemma}

% [AMENDED]
% Owner: C
\begin{lemma}\label{l:Ostruct}
For \(\mathcal I\) a union of vector groups and
\(\mathcal O=\mathcal O(\widehat F_c^{\mathcal I})\),
the matrix in the standard tensor basis is real symmetric.
It commutes with \(U^{\otimes k}\) for every \(U\in\mathrm{SU}(2)\)
and with \(S_z,S_+,S_-\).
For a retained \(\mathrm U\)-corner pair, if \(\Pi\) projects that
pair onto \(\operatorname{Sym}^2\mathbb C^2\), then
\(\mathcal O=\Pi\mathcal O\Pi\).
For a retained \(\mathrm D\)-corner pair, \(\mathcal O\) commutes
with exchange of its two factors.
\end{lemma}

% [AMENDED]
% Owner: C
\begin{definition}\label{d:exchange}
Let \(V\) be the tensor product of \(\mathbb C^2\) for singles,
\(\operatorname{Sym}^2\mathbb C^2\) for retained \(\mathrm U\)-pairs,
and \(\mathbb C^2\otimes\mathbb C^2\) for retained \(\mathrm D\)-pairs.
For each assignment \(\epsilon\) of exchange signs to the
\(\mathrm D\)-pairs, let \(V_\epsilon\) replace each of those
pair factors by its symmetric or exterior square.
If there are \(t\) singles, let \(\kappa=0\) for even \(t\),
\(\kappa=1/2\) for odd \(t\), and
\(V_\epsilon^\circ=\ker(S_z-\kappa I)|_{V_\epsilon}\).
The standard sector basis uses \(e_\pm\) for singles and
\(e_+\otimes e_+\), \(e_+\otimes e_-+e_-\otimes e_+\),
\(e_-\otimes e_-\) for symmetric pairs, with
\(e_+\otimes e_--e_-\otimes e_+\) for antisymmetric pairs.
\end{definition}

% [AMENDED]
% Owner: C
\begin{lemma}\label{l:sector}
Let \(V'\) be any tensor sector formed from singles and
symmetric or exterior pair factors as above.
If a Hermitian \(\mathcal O'\) on \(V'\) commutes with
\(S_z,S_+,S_-\), then it preserves the central-weight space
\(V'^\circ\), and its restriction there has the same set of
eigenvalues as \(\mathcal O'\).
In particular \(\mathcal O'\) is positive definite if and only if
that restriction is positive definite.
\end{lemma}

% [AMENDED]
% Owner: C
\begin{lemma}\label{l:gram-residual}
Let \(M\) be a symmetric integer matrix of size \(n\).
Let \(g,d,S\) be positive integers and \(L\) an integer
lower-triangular \(n\times n\) matrix; no invertibility is required.
Define the symmetric integer residual
\[
 E=dS^2M-gLL^{\mathsf T}.
\]
If, for every row,
\[
 E_{ii}>0,\qquad E_{ii}>\sum_{j\ne i}|E_{ij}|,
\]
then \(E\) and \(M\) are positive definite over \(\mathbb R\)
and, as Hermitian matrices, over \(\mathbb C\).
The case \(L=0\), \(g=d=S=1\), is the direct strict
diagonal-dominance certificate.
All hypotheses can be checked with integer sums, products,
absolute values and comparisons, without division or determinants.
\end{lemma}

% [AMENDED]
% Owner: C
\begin{lemma}\label{l:integercheck}
For \(c=p/q\in(0,1)\) rational with positive integers \(p,q\),
choose a positive integer \(D_0\) clearing the coefficients of
\(3D_0\widehat Z_{\rm cur}\) and \(5D_0\widehat Z_{\rm next}\).
In the integer-coordinate central-sector basis, let \(A_\epsilon\)
and \(A'_\epsilon\) be the Hermitian-form matrices of their
inverted Pauli operators. These are symmetric integer matrices.
The form of \(\mathcal O(\widehat F_c^{\mathcal I})\) on the sector
is positive definite if and only if
\[
 M_\epsilon=5qA_\epsilon-3pA'_\epsilon
\]
is positive definite. The primary verification uses the integer
Gram/residual witness of Lemma~\ref{l:gram-residual}.
As an optional fallback only, positivity is equivalent to positivity of all its
leading principal minors. More explicitly, for a size-\(t\) matrix
set \(A^{(0)}=M_\epsilon\), \(p_{-1}=1\), and for \(0\le j<t\)
set \(p_j=A^{(j)}_{jj}\) and update the trailing entries by
\[
 A^{(j+1)}_{il}
    =\frac{p_jA^{(j)}_{il}-A^{(j)}_{ij}A^{(j)}_{jl}}{p_{j-1}}
       \quad(i,l>j).
\]
As long as preceding pivots are nonzero, each division is exact
in the integers and \(p_j\) is the leading principal minor of
order \(j+1\). Acceptance requires every division exact and every
pivot positive.
\end{lemma}

% [AMENDED]
% Owner: C
\begin{proposition}\label{p:certificate}
Let \(c=p/q\in(0,1)\), with \(p,q\in\mathbb Z_{>0}\).
Choose either allowed representation, an inversion which is a
union of groups, and a justified positive denominator \(D_0\).
For every supported central exchange sector reconstruct
\(M_\epsilon=5qA_\epsilon-3pA'_\epsilon\) as in
Lemma~\ref{l:integercheck}.
Suppose that for each sector there are positive integers
\(g_\epsilon,d_\epsilon,S_\epsilon\) and an integer lower-triangular
matrix \(L_\epsilon\) of its dimension such that
\[
 E_\epsilon=d_\epsilon S_\epsilon^2M_\epsilon
                  -g_\epsilon L_\epsilon L_\epsilon^{\mathsf T},
 \qquad
 (E_\epsilon)_{ii}>
          \sum_{j\ne i}|(E_\epsilon)_{ij}|\quad\text{for every }i.
\]
Then \(F_c(m)\ge0\) for every \(m\in(\mathbb S^2)^{20}\).
The operator may vanish on the discarded antisymmetric
\(\mathrm U\)-pair subspaces; positive definiteness is required
only on the supported exchange sectors.
\end{proposition}


\end{document}
